% Aula 02 --- a Ridge como moda a posteriori sob priori gaussiana.
Já mencionamos que o Ridge tem uma leitura bayesiana: ele é a moda a
posteriori sob uma priori gaussiana. A afirmação é mais do que curiosidade ---
ela dá um significado concreto ao $\lambda$, que até aqui era só um hiperparâmetro.

Considere o modelo
\[
  \bm y \mid \bbeta \;\sim\; N\big(\mathbb{X}\bbeta,\ \sigma^2\mathbb{I}\big),
  \qquad
  \bbeta \;\sim\; N\big(\bm 0,\ \sigma_\beta^2\mathbb{I}\big),
\]
com $\sigma^2$ e $\sigma_\beta^2$ conhecidos.
\begin{itens}
  \item \pts{1,6} Escreva a densidade a posteriori de $\bbeta$ a menos de
        constante multiplicativa e mostre que maximizá-la equivale a minimizar
        \[
          \frac{1}{2\sigma^2}\norm{\bm y-\mathbb{X}\bbeta}^2
          + \frac{1}{2\sigma_\beta^2}\norm{\bbeta}^2 .
        \]
        O caminho é o de sempre: passar ao logaritmo transforma produto em soma, e
        trocar o sinal transforma maximizar em minimizar.
  \item \pts{1,0} Compare o que você obteve com a definição do Ridge da aula e
        identifique $\lambda$ em função de $\sigma^2$ e $\sigma_\beta^2$. Cuidado
        com o fator que precisa ser eliminado antes da comparação.
  \item \pts{1,4} Agora a parte interessante: o que a priori está \emph{afirmando}
        quando $\sigma_\beta^2\to\infty$, e quando $\sigma_\beta^2\to 0$? Ligue
        cada resposta ao comportamento do Ridge em $\lambda\to 0$ e
        $\lambda\to\infty$.
\end{itens}

\begin{solucao}
\textbf{(a)} Pelo teorema de Bayes, a posteriori é proporcional ao produto da
verossimilhança pela priori:
\[
  \pi(\bbeta\mid\bm y)\;\propto\; f(\bm y\mid\bbeta)\,\pi(\bbeta)
   \;\propto\;
   \exp\Big\{-\frac{1}{2\sigma^2}\norm{\bm y-\mathbb{X}\bbeta}^2\Big\}
   \cdot
   \exp\Big\{-\frac{1}{2\sigma_\beta^2}\norm{\bbeta}^2\Big\}.
\]
As constantes de normalização das duas gaussianas não dependem de $\bbeta$ e por
isso podem ser absorvidas no $\propto$.

Maximizar uma função positiva é o mesmo que maximizar seu logaritmo, e o
logaritmo de um produto de exponenciais é a soma dos expoentes:
\[
  \log\pi(\bbeta\mid\bm y)
   = \text{const} - \frac{1}{2\sigma^2}\norm{\bm y-\mathbb{X}\bbeta}^2
     - \frac{1}{2\sigma_\beta^2}\norm{\bbeta}^2 .
\]
Maximizar isso é minimizar o simétrico dos dois termos que dependem de $\bbeta$,
que é exatamente o objetivo do enunciado.

\smallskip
\textbf{(b)} O objetivo de (a) tem um $1/(2\sigma^2)$ na frente do primeiro termo,
e a definição do Ridge não tem fator nenhum ali. Multiplique tudo por
$2\sigma^2>0$ --- o que não altera o minimizador, por ser constante positiva:
\[
  \norm{\bm y-\mathbb{X}\bbeta}^2 + \frac{\sigma^2}{\sigma_\beta^2}\norm{\bbeta}^2 .
\]
Agora a comparação é imediata, e
\[
  \boxed{\;\lambda=\frac{\sigma^2}{\sigma_\beta^2}\;}
\]

\smallskip
\textbf{(c)} Quando $\sigma_\beta^2\to\infty$, a priori fica \textbf{difusa}: ela
espalha massa por toda parte e não expressa preferência alguma por coeficientes
pequenos. Em outras palavras, ela não afirma nada. Correspondentemente
$\lambda\to 0$ e o Ridge degenera no MQO: os dados decidem sozinhos, com viés
baixo e variância alta.

Quando $\sigma_\beta^2\to 0$, a priori fica \textbf{concentrada em zero}: ela
afirma, antes de ver dado nenhum, que os coeficientes são praticamente nulos.
Correspondentemente $\lambda\to\infty$ e $\widehat\bbeta\to\bm 0$: viés máximo,
variância mínima.

Entre os dois extremos, a fórmula $\lambda=\sigma^2/\sigma_\beta^2$ diz uma coisa
bem razoável: regulariza-se mais quando os \emph{dados} são ruidosos
($\sigma^2$ grande) ou quando a \emph{crença prévia} é firme ($\sigma_\beta^2$
pequeno). O $\lambda$ é a razão entre duas incertezas, e escolhê-lo por validação
cruzada é, nessa leitura, deixar os dados escolherem a priori que melhor prediz.
\end{solucao}
